\begin{pastexam}[label={px:ncku-114-cce-ds-01}]{成大電通所 114 學年度 單選題第 1 題}
（單選；每題 5 分，答錯倒扣 2 分）

A is a two dimension array, the locations of A(3, 2) and A(2, 3) are 1115 and 1207, respectively. Assume
that each element occupies two addresses. Please give the location of A(5, 10).
\begin{enumerate}[label=\alph*),nosep]
\item 1841
\item 1851
\item 1861
\item 1871
\item None of the above.
\end{enumerate}
\end{pastexam}

\begin{pastsoln}
答案：d) 1871。

題目沒說是以列為主（row-major）還是以行為主（column-major），也沒給索引的下界。 下界不必知道，只要用位址的差；
存放方式用兩個已知位址判斷（資結書第~4.1~節）。 每個元素佔 2 個位址。

\textbf{以列為主}：設每列有 $C$ 個元素，$\mathrm{loc}(i, j) = \alpha + 2\bigl[(i - l_1)C + (j - l_2)\bigr]$，所以
\[
\mathrm{loc}(3, 2) - \mathrm{loc}(2, 3) = 2\bigl[(3 - 2)C + (2 - 3)\bigr] = 2(C - 1) = 1115 - 1207 = -92,
\]
得 $C = -45$，不可能。 所以不是以列為主。

\textbf{以行為主}：設每行有 $R$ 個元素，$\mathrm{loc}(i, j) = \alpha + 2\bigl[(j - l_2)R + (i - l_1)\bigr]$，
\[
\mathrm{loc}(2, 3) - \mathrm{loc}(3, 2) = 2\bigl[(3 - 2)R + (2 - 3)\bigr] = 2(R - 1) = 1207 - 1115 = 92,
\]
得 $R = 47$。 於是
\[
\mathrm{loc}(5, 10) = \mathrm{loc}(3, 2) + 2\bigl[(10 - 2) \cdot 47 + (5 - 3)\bigr] = 1115 + 2 \times 378 = 1115 + 756 = 1871.
\]
選 d)。 （可以再檢查：$\mathrm{loc}(2, 3) + 2[(10 - 3) \cdot 47 + (5 - 2)] = 1207 + 2 \times 332 = 1871$，一致。）
\end{pastsoln}
